\def\topic{Morse-Kelley class comprehension, separation, and Dunaev tagged ordered pairs}
\ifdefined\MyHeaderIncluded
\else
\def\qwertyflag{true}

\ifdefined\MyHeaderIncluded

\else
  \newcommand{\foobar}{FooBar}
  \def\MyHeaderIncluded{true}
% if is till the end of the file

\documentclass[10pt,a4paper]{article}
\usepackage[utf8]{inputenc}
%\usepackage[german]{babel}
%\usepackage[cyrillic]{babel}
\usepackage[T1]{fontenc}
%\usepackage[T2A]{fontenc}
%\usepackage[english, russian]{babel}
\usepackage{amsmath}
\usepackage{amsfonts}
\usepackage{amssymb}
\usepackage{amsthm}
\usepackage{mathtools} %coloneqq

\usepackage{bussproofs} %

%\renewcommand{\thesubsection}{\arabic{subsection}}
%\newtheorem{thm}{Theorem}[section]
%\renewcommand{\thethm}{\arabic{subsection}.\arabic{thm}}

\usepackage{thmtools}
\usepackage{hyperref}
\hypersetup{
    colorlinks=true,
    linkcolor=blue,
    filecolor=magenta,      
    urlcolor=cyan,
    pdftitle={Overleaf Example},
    %pdfpagemode=FullScreen,
    }
%\usepackage{stix-type1}%\turnediota
\theoremstyle{definition}

\newtheorem{constructor}{Constructor}[section]
\newtheorem{iaxiom}{Axiom}
\newtheorem{axiom}{Axiom}
\newtheorem*{axiom*}{Axiom}
\newtheorem{definition}{Definition}[section]
\newtheorem*{definition*}{Definition}
\newtheorem{remark}{Remark}[section]
\newtheorem*{remark*}{Remark}
\newtheorem{recall}{Recall}[section]
\newtheorem*{recall*}{Recall}

\newtheorem{lemma}{Lemma}[section]
\newtheorem{theorem}{Theorem}[section]
\newtheorem*{theorem*}{Theorem}
\newtheorem{note}{Note}[section]
\usepackage{etoolbox}

%\newcommand*{\@Dashv}[2]{%
%  \reflectbox{$\m@th#1#2$}%
%}
\newcommand*{\separation}[3]{% 1)binder, 2)set, 3)condition
  \{#1\in #2: #3\}
}

% https://www.freecodecamp.org/news/simple-regex-tricks-for-beginners-3acb3fa257cb/
%  \\mathcal\{T\} |---> \topology
%  \\mathcal\{([^\}]*)\}   |--->  $1
\newcommand{\AND}{\ \ \&\ \ }
\newcommand{\any}{\mathrm{\ any}}
\newcommand{\borel}{\mathcal{B}}
\newcommand{\can}{\mathrm{can}}
\newcommand{\canonize}{\mathrm{canonize}}
\newcommand{\closure}{\mbox{cl}}
\newcommand{\complex}{\mathbb{C}}
\newcommand{\converges}{\longrightarrow}
\newcommand{\countable}{\mbox{\ countable}}
\newcommand{\ldefined}{\mbox{\ defined}}
\newcommand{\dom}{\mbox{dom}}
\newcommand{\Dom}{\mbox{Dom}}
\newcommand{\Done}{\noindent \textbf{Done.}\\}
\newcommand{\lexists}{\mbox{\ exists}}
\newcommand{\lforall}{\mbox{\ forall}}
\newcommand{\lunique}{\mbox{\ unique}}
\newcommand{\family}{\mathrm{\ family}}
\newcommand{\Fm}{\mathrm{Fm}}
\newcommand{\FuncSymb}{\mathrm{FuncSymb}}
\newcommand{\function}{\mathrm{\ function}}
\newcommand{\functional}{\mbox{\ functional}}
\newcommand{\IFF}{\ \Longleftrightarrow\ }
\renewcommand{\iff}{\ \longleftrightarrow\ }
\newcommand{\IMPLIES}{\ \Longrightarrow\ }
\renewcommand{\implies}{\ \longrightarrow\ }
\newcommand{\Inf}{\mathrm{Inf}}
\newcommand{\injective}{\mbox{\ injective}}
\newcommand{\injection}{\mbox{\ injection}}
\newcommand{\interior}{\mbox{int}}
\newcommand{\intersection}{\bigcap}
\newcommand{\isempty}{\mathrm{\ empty}}
\newcommand{\LB}{\mathrm{LB}}
\newcommand{\nghbset}{\mathcal{N}}
\newcommand{\nonempty}{\mathrm{\ nonempty}}
\newcommand{\Ord}{\mathrm{Ord}}
\newcommand{\power}{\mathcal{P}}
\newcommand{\pr}{\mathrm{pr}}
\newcommand{\fst}{\mathrm{fst}}
\newcommand{\snd}{\mathrm{snd}}
\newcommand{\PredSymb}{\mathrm{PredSymb}}
\newcommand{\Proof}{\noindent \textbf{Proof.}\\}
\newcommand{\ran}{\mathrm{ran}}
\newcommand{\Ran}{\mathrm{Ran}}
\newcommand{\real}{\mathbb{R}}
\newcommand{\relation}{\mathrm{\ relation}}
\newcommand{\set}{\mathrm{\ set}}
\newcommand{\singleton}{\mathrm{\ singleton}}
\newcommand{\semantics}{\mathrm{semantics}}
\newcommand{\semialgebra}{\mathcal{S}}
\newcommand{\Sup}{\mathrm{Sup}}
\newcommand{\surjective}{\mbox{\ surjective\ on\ }}
\newcommand{\Term}{\mathrm{Term}}
\newcommand{\Type}{\mathrm{Type}}
\newcommand{\topology}{\mathcal{T}}
\newcommand{\total}{\mbox{\ total\ on\ }}
\newcommand{\transitive}{\mathrm{\ transitive}}
\newcommand{\union}{\bigcup}
\newcommand{\lif}{\mathrm{if}}
\newcommand{\lthen}{\mathrm{then}}
\newcommand{\lelse}{\mathrm{else}}
\newcommand{\lexception}{\mathrm{exception}}
\newcommand{\lundef}{\mathrm{undef}}
\newcommand{\univocal}{\mathrm{\ univocal}}
\newcommand{\univocalat}{\mathrm{\ univocal\ at\ }}
\newcommand{\UB}{\mathrm{UB}}
\newcommand{\Var}{\mathrm{Var}}

\newcommand{\notvarphi}{\phi}
\renewcommand{\phi}{\varphi}

\newcommand{\comment}[1]{\quad\mbox{(#1)}}
% local commands
%\preto\section{\undef\blipo} %etoolbox
\preto\subsection{\undef\blipo} %etoolbox
% https://tex.stackexchange.com/questions/105684/limit-macro-newcommand-scope-to-section
\usepackage{graphicx}
\usepackage{authblk}

\setlength\parindent{0pt} % whole doc of \noindent

%\usepackage{mathgl}

\fi % guard
\ifdefined\topic
\else
\def\topic{No Titel}
\fi
\title{\topic}
\author[1]{Georgy Dunaev}
\begin{document}
\maketitle
\fi
\section{Morse-Kelley Class Separation and Dunaev Tagged Ordered Pairs}
\label{sec:mk_classes_audit_summary}

This standalone unit summarizes the core Morse-Kelley (MK) class-theoretic
principles established in \texttt{100\_class\_theory.tex}, \texttt{010\_ordered\_pairs.tex},
and \texttt{070\_ordered\_pair.tex}, focusing on class comprehension, the
Separation Principle for sets intersected with arbitrary classes, and the
mandate to use Dunaev Tagged Ordered Pairs rather than Kuratowski pairs for
general classes.

\subsection{Universal Class, Sethood, and Class Comprehension}

Let $V$ denote the universal class of all sets:
\[
V \eqdef \cc{x}{\top}.
\]
A class $A$ is a \emph{set} if $A \in V$, and a \emph{proper class} if $A \notin V$.
For every formula $\phi(x)$ in the language of Morse-Kelley class theory (in
which bound variables may range over arbitrary classes), the Class Comprehension
Axiom Schema provides a class $\cc{x}{\phi(x)}$ satisfying
\[
\forall z,\ z \in \cc{x}{\phi(x)} \iff (z \in V) \land \phi(z).
\]
In particular, the Russell class $R \eqdef \cc{x}{x \notin x}$ satisfies
$R \in R \iff (R \in V) \land (R \notin R)$, which yields $R \notin V$; hence $R$
is a proper class.

\subsection{Separation Principle (Aussonderung)}

In Morse-Kelley class theory, intersecting a set $a \in V$ with an arbitrary
class $P$ (which may be defined by an impredicative formula quantifying over all
classes) yields the subclass
\[
a \cap P = \cc{x}{x \in a \land x \in P}.
\]
The Separation Principle states that every subclass of a set is itself a set:
\[
\forall a,\ \forall P,\ (a \in V) \implies (a \cap P \in V).
\]
As shown in Theorem~\ref{thm:two_object_model} of \texttt{100\_class\_theory.tex},
class comprehension and extensionality alone determine the extension of $a \cap P$
as a class, whereas its sethood $a \cap P \in V$ is an independent set-existence
axiom (derivable in full MK from the Limitation of Size or Replacement plus the
empty-set axiom).

\subsection{Adherence to Dunaev Tagged Ordered Pairs for Classes}

For sets $a, b \in V$, the Kuratowski pair
\[
\kuratowskipair{a}{b} \eqdef \{\{a\}, \{a, b\}\}
\]
relies on $\{a\} \in V$ and $\{a, b\} \in V$. For a proper class $A \notin V$,
however, $\{A\} = \cc{x}{x = A} = \emptyclass$, so $\kuratowskipair{A}{B}$
collapses when $A$ or $B$ is a proper class. Accordingly, Kuratowski pairs must
never be used for proper classes.

To form ordered pairs of arbitrary classes $A$ and $B$ without sethood
restrictions, the development uses \emph{Dunaev Tagged Ordered Pairs}:
\[
\pair{A}{B} \eqdef \cc{z}{\exists x,\ (z = \kuratowskipair{\emptyclass}{x} \land x \in A) \lor (z = \kuratowskipair{\{\emptyclass\}}{x} \land x \in B)},
\]
with projections
\[
\Fst(P) \eqdef \cc{x}{\kuratowskipair{\emptyclass}{x} \in P},
\qquad
\Snd(P) \eqdef \cc{x}{\kuratowskipair{\{\emptyclass\}}{x} \in P}.
\]
Because every element $x \in A$ or $x \in B$ is a set in $V$ and the tags
$\emptyclass, \{\emptyclass\}$ are distinct sets in $V$, the tagged elements
$\kuratowskipair{\emptyclass}{x}$ and $\kuratowskipair{\{\emptyclass\}}{x}$ are
always sets in $V$. Consequently, $\Fst(\pair{A}{B}) = A$, $\Snd(\pair{A}{B}) = B$,
and $\pair{A}{B} = \pair{C}{D} \iff (A = C \land B = D)$ hold unconditionally for
all classes $A, B, C, D$, whether sets or proper classes.
\ifdefined\qwertyflag
\undef\qwertyflag
\end{document}
\fi
